\section{Discussion}
\label{sec:discussion}

\subsection{Why a local object and a global object come apart}
$\chieq$ is a derivative evaluated at a single equilibrium; $\EPR$ is a functional of a flux
field over the entire profile space. Near $\alpha=0$ the potential component modulates both,
which is why they track: the potential is a scalar field whose gradient governs the local
response and whose level sets govern the stationary measure. As the potential component
vanishes, nothing remains to force a local object and a global object into agreement, and
they stop agreeing. This is not a statement peculiar to games. It predicts that in any
system where a local susceptibility and a global flux functional are both used as proxies for
irreversibility, the two will diverge precisely where the underlying potential structure
does---and it says where to look.

\subsection{What $\R$ is not}
$\R$ is not a measure of dissipation, and this paper's main practical warning is against
reading it as one. It is also not a fraction: it is a ratio of Frobenius norms, unbounded
above, and values greater than $1$ occur (matching pennies reads $1.2$). Its \emph{zero
test} is $\lambda$-free---$\R=0$ if and only if the normalised game is potential, at every
$\lambda$---but its \emph{magnitude} scales with $\lambda$, and an earlier version of this
programme's own claim that $\R$ was ``$\lambda$-free'' was corrected on exactly this point.
Any reported magnitude must carry the $\lambda$ it was read at.

\subsection{Consequences for practice}
The plane has four quadrants and they differ in what a decision-maker should do.
Reciprocal-and-quiet systems admit comparative statics and static rival models.
Reciprocal-but-circulating systems are being driven from outside, so timing matters and
structure does not. Asymmetric-but-quiet systems have a structural leader, so pass-through
asymmetry is the exploitable object and there is no cycle to time. Systems that are both are
where naive optimisation against a static rival model is worst. If distance from equilibrium
were a scalar, the two off-diagonal quadrants would be empty; the two real systems measured
here are not on the same diagonal.

\subsection{Theoretical connections}
The zero of the local axis is the strategic instance of Onsager reciprocity, and the zero of
the global axis is detailed balance; that they coincide in potential games is the exact
content of the Gibbs correspondence. Away from it, the structure resembles the failure of
the fluctuation--dissipation theorem in a non-equilibrium steady state---with the important
difference that here the quantity governing the failure, $\alpha$, is itself computable from
the game, so the breakdown can be parameterised rather than merely observed.

\subsection{Limitations}
The finite-size limitation has been tested rather than merely acknowledged, and the outcome
is \textbf{indeterminate}. Sweeping $m\in\{3,4,5,6\}$ under criteria registered in advance
(Section~\ref{sec:robustness}), two of the three conditions that would have killed the claim
fired; the third---a demand that the collapse lose more than half its size---did not, and the
gap ratio deciding it, $0.585$, carries a bootstrap interval that crosses its own $0.50$ bar.
Under the \emph{earlier} of this unit's two registrations, which adjudicates the raw recovery
$\Delta\rho_{\mathrm{hi}}=+0.364$ against bars of $0.40$ and $0.20$, neither branch resolves.
We therefore report indeterminate and rest the claim on the ceiling condition instead. What
the sweep does establish without qualification is that the collapse itself, its ceiling, and
the $\alpha$ at which it begins are present at every action count tested.

\textbf{A criterion substitution occurred inside this unit, and it is reported here rather
than left to be discovered.} The kill-shot was first registered as
\texttt{plane\_finite\_size.yaml} at 23:02:25 on 2026-08-12; an experiment implementing it
was written at 23:04:39, executed at 23:04:49, and run again at 23:17:47. At 23:26:36---after
both runs---a second configuration, \texttt{plane\_robustness.yaml}, was committed for the
same unit and the same kill-shot with a different deciding statistic, and the change ran in
the direction of survival. Each file individually satisfies the discipline this programme
requires (committed before the code implementing it existed; never amended after a result
existed); the \emph{unit} does not, because substituting a criterion by writing a new file
leaves no diff to catch. Both registrations, both experiments and both resolved
configurations are committed, so the property is checkable rather than asserted. Two
consequences are carried in the text above: the reported verdict is the earlier
registration's, and the deciding statistic of the later one now carries the interval the
registration failed to demand of it---an interval that does not certify the branch it was
read as certifying. The general lesson is that ``the config was committed before the code''
is not a sufficient guarantee; the binding property is that a unit has \emph{one} live
criterion for one question, and that any replacement is registered as a replacement.

\textbf{The control arms are unpaired, which is a design defect rather than a caveat.} The
payoff-scale control D1 and the precision control D2 draw from different seed offsets than
the main arm and use half the games per cell, so every arm-to-arm comparison in
Section~\ref{sec:robustness} carries between-sample noise that a paired design---same games,
same order, only the controlled quantity varying---would cancel. The solver check H6 is
paired, which is exactly why its bar can be strict. D1 and D2 should be re-run paired, and
D1 additionally with $\bar\lambda$ fixed by construction rather than bracketed: as run, D1
raises $\bar\lambda$ with $m$ where the main arm lowers it, and interpolating between the two
arms to matched $\bar\lambda$ leaves most of the observed $m$-effect standing. D2 is a
further signal that the gap-ratio criterion is not resolvable at this sample size: at the
precision $\lambda=1.5$ it returns a gap ratio of $0.529$, a $2.9$-point margin against the
same $0.50$ bar on $100$ games per cell.

The second limitation is dimension in the other index: \textbf{everything here is at
$N=2$}. Nothing in this paper speaks to $N$-scaling, and it is the obvious next kill-shot---
the same design, run over $N$ rather than $m$, with the same three-branch criterion. The
binding cost is the dense Glauber generator over $\prod_i m_i$ profiles, not a correctness
limit; $m$ stops at $6$ (36 joint states) for the same reason.

A third limitation is the sign, and the $m$-sweep sharpened it. The reversal is one seed
family's realisation of an effect whose sign varies with $\lambda$ and, at fixed Frobenius
payoff scale, with $m$: at $m=6$ the terminal association is statistically indistinguishable
from zero. The paper's claim is therefore decorrelation, not anti-correlation---which is all
independence requires---and any use of the reversal must carry the $(m,\lambda)$ it was read
at. Table~\ref{tab:scoreboard} records H3b as not supported and the reversal half of H5 as
not general, rather than presenting either as established.

A fourth limitation is generic to the generator and is a warning to anyone extending this
work: because the family unit-normalises each Hodge component and fixes the Frobenius norm,
per-entry payoff RMS falls like $1/m$, so a larger action set is also a colder game. \textbf{Any
future sweep in $m$ on this family must control per-entry payoff RMS, or it will confound
temperature with dimension.} The registered control (Table~\ref{tab:d1}) is how the confound
was caught here, and it is also why the $m$-trend cannot be attributed to temperature
outright: the control brackets $\bar\lambda$ rather than matching it, and at matched
$\bar\lambda$ a residual $m$-effect of roughly $0.22$ of the observed $0.364$ remains. The
same caution applies to comparing $\R$ or $\EPR$ magnitudes across action
counts at fixed Frobenius scale.

Solver dependence is no longer an open limitation, though the check that closed it had less
power than its margin suggests. H6---a structurally different, last-iterate-convergent method
and a $100\times$ tighter tolerance, on the identical seed family, criterion
$|\Delta\rho_S|<0.05$---returned $0.00000$ in all four comparisons. Because the two arms move
$\R$ by only $\sim\!10^{-12}$ and $\rho_S$ is rank-based, that outcome was forced rather than
earned; what the check genuinely establishes is the diagnostic that the solver contributes
$\sim\!10^{-12}$ to $\R$ and leaves its rank ordering bit-identical, which is enough to
exclude finite-precision path dependence. What remains untested is a bias shared by both
methods, which no comparison of two converging solvers can exclude; exact homotopy
continuation would be the next distinct probe.

A fifth limitation is that no system has yet been measured on \emph{both} coordinates. The
retail panel has $\R$ without $\EPR$; the electricity market has $\EPR$ without $\R$. Until
one system carries both, the plane is demonstrated on synthetic families and anchored, not
populated, by real data. This is a real limitation of the present evidence rather than
evidence of generality.

A sixth limitation concerns the empirical anchors themselves. The retail panel is a single
chain in one metropolitan area over 1989--1997, and its cross-``firm'' interpretation is a
modelling assumption, not a data feature; the near-symmetric reading is consistent with a
single retailer optimising a category, which is what was predicted, but it is not evidence
about rival firms.

Finally, prior art. The nearest live work~\cite{legacci2024} was assessed as orthogonal, but
that assessment is inherited from an internal literature audit and has not been independently
re-run for this paper. It is registered as an open check.

\subsection{Broader relevance}
The two coordinates require no knowledge of payoffs, which is what makes them usable outside
game theory: any system that emits a sequence of joint actions and permits a perturbation
experiment can be located in the plane. Congestion networks, wholesale electricity markets,
retail pricing panels, and multi-agent learning systems all satisfy that description, and
the first three appear here.
